\chapter{Emission Model Derivations}
\label{appendix:emission_updates}

This appendix presents the theoretical derivations of the emission-model updates used in the proposed \ac{DBN}. 
The derivation follows the standard \ac{EM} formulation for latent-variable models
\cite{Dempster1977,pml1Book}. Exact implementation details and training-specific choices are described in the methodology chapter.

Let \(O_{1:T}=(O_1,\dots,O_T)\) denote the observed sequence, and let
\(S_{1:T}\) and \(A_{1:T}\) denote the corresponding latent style and action
state sequences. The model parameters are collected in \(\theta\).
Each observation at time step \(t\) is decomposed into a continuous part and a binary part:
\begin{equation}
O_t = (x_t, b_t)
\end{equation}
with
\begin{align}
x_t &\in \mathbb{R}^{D_c}, \\
b_t &\in \{0,1\}^{B}
\end{align}


In the E-step of the \ac{EM} algorithm, the current parameter estimate
\(\theta^{\mathrm{old}}\) is used to compute posterior probabilities over the latent variables \cite{rabiner1989tutorial,murphy2002dynamic,pml1Book}.

The main posterior marginal used in the emission updates is:
\begin{equation}
\gamma_t(s,a)
=
p(S_t=s,A_t=a \mid O_{1:T},\theta^{\mathrm{old}})
\label{eq:gamma_joint_appendix}
\end{equation}
This quantity is the posterior responsibility assigned to the latent pair \((s,a)\) at time step \(t\). 
For each fixed \(t\), the values \(\{\gamma_t(s,a)\}_{s,a}\) form a normalized joint posterior distribution:
\begin{equation}
\sum_{s,a}\gamma_t(s,a)=1
\end{equation}

The transition-update derivations also use the posterior transition marginals:
\begin{align}
\xi_t^{S}(s',s)
&=
p(S_{t-1}=s',S_t=s \mid O_{1:T},\theta^{\mathrm{old}}),
\qquad t=2,\dots,T,
\label{eq:xi_style_appendix}
\\
\xi_t^{A}(s,a',a)
&=
p(S_t=s,A_{t-1}=a',A_t=a \mid O_{1:T},\theta^{\mathrm{old}}),
\qquad t=2,\dots,T
\label{eq:xi_action_appendix}
\end{align}


\subsection{EM Objective}

The goal of learning is to maximize the observed-data log-likelihood:
\begin{equation}
\log p(O_{1:T}\mid \theta)
=
\log \sum_{S_{1:T},A_{1:T}}
p(O_{1:T},S_{1:T},A_{1:T}\mid \theta)
\label{eq:obs_loglik}
\end{equation}
Direct maximization of \eqref{eq:obs_loglik} is difficult because the latent state
sequences are not observed and must be summed out \cite{Dempster1977,pml1Book}.

In the \ac{EM} algorithm, this difficulty is addressed by introducing an auxiliary distribution \(q(S_{1:T},A_{1:T})\) over the latent trajectories, with:
\begin{equation}
q(S_{1:T},A_{1:T}) \ge 0, 
\qquad 
\sum_{S_{1:T},A_{1:T}} q(S_{1:T},A_{1:T}) = 1
\end{equation}
Using this distribution, the log-likelihood can be rewritten as:
\begin{align}
\log p(O_{1:T}\mid \theta)
&=
\log \sum_{S_{1:T},A_{1:T}}
q(S_{1:T},A_{1:T})
\frac{p(O_{1:T},S_{1:T},A_{1:T}\mid \theta)}
{q(S_{1:T},A_{1:T})}
\label{eq:em_rewrite}
\end{align}

Jensen's inequality states that for a convex function \(f\),
\begin{equation}
f\!\left(\sum_i q_i x_i\right)
\le
\sum_i q_i f(x_i),
\qquad
q_i \ge 0,
\quad
\sum_i q_i = 1.
\label{eq:jensen_convex}
\end{equation}
Equivalently, if \(f\) is concave, the inequality direction is reversed. Since
\(\log(\cdot)\) is concave, applying Jensen's inequality to \eqref{eq:em_rewrite} gives:
\begin{align}
\log p(O_{1:T}\mid \theta)
&\ge
\sum_{S_{1:T},A_{1:T}}
q(S_{1:T},A_{1:T})
\log
\frac{p(O_{1:T},S_{1:T},A_{1:T}\mid \theta)}
{q(S_{1:T},A_{1:T})}
\label{eq:em_lower_bound}
\end{align}

Expanding the logarithm yields:
\begin{align}
\log p(O_{1:T}\mid \theta)
&\ge
\sum_{S_{1:T},A_{1:T}}
q(S_{1:T},A_{1:T})
\log p(O_{1:T},S_{1:T},A_{1:T}\mid \theta)
\nonumber\\
&\quad-
\sum_{S_{1:T},A_{1:T}}
q(S_{1:T},A_{1:T})\log q(S_{1:T},A_{1:T})
\label{eq:em_bound_expanded}
\end{align}
The second term is the entropy of \(q\). For fixed \(q\), it does not depend on
\(\theta\). Therefore, in the M-step it is sufficient to maximize only the first term \cite{Dempster1977,pml1Book}.


In \ac{EM}, the auxiliary distribution is chosen as the posterior under the current
parameter estimate:
\begin{equation}
q(S_{1:T},A_{1:T})
=
p(S_{1:T},A_{1:T}\mid O_{1:T},\theta^{\mathrm{old}})
\label{eq:q_posterior_choice}
\end{equation}
Substituting this choice into the lower bound gives the auxiliary function:
\begin{equation}
Q(\theta,\theta^{\mathrm{old}})
=
\mathbb{E}_{p(S_{1:T},A_{1:T}\mid O_{1:T},\theta^{\mathrm{old}})}
\left[
\log p(O_{1:T},S_{1:T},A_{1:T}\mid \theta)
\right]
\label{eq:q_function_def}
\end{equation}

Thus, the E-step computes the posterior over the latent trajectories under
\(\theta^{\mathrm{old}}\), and the M-step updates \(\theta\) by maximizing the
expected complete-data log-likelihood \(Q(\theta,\theta^{\mathrm{old}})\)
\cite{Dempster1977,rabiner1989tutorial,murphy2002dynamic,pml1Book}.




Using the factorization of the proposed \ac{DBN}, the complete-data distribution can be written as:
\begin{align}
p(O_{1:T},S_{1:T},A_{1:T}\mid \theta)
&=
p(S_1\mid \theta)\,
p(A_1\mid S_1;\theta)\,
\prod_{t=2}^{T} p(S_t\mid S_{t-1};\theta)
\nonumber\\
&\quad\times
\prod_{t=2}^{T} p(A_t\mid A_{t-1},S_t;\theta)\,
\prod_{t=1}^{T} p(O_t\mid S_t,A_t;\theta)
\label{eq:complete_data_factorization}
\end{align}


Taking the logarithm converts the product into a sum:
\begin{align}
\log p(O_{1:T},S_{1:T},A_{1:T}\mid \theta)
&=
\log p(S_1\mid \theta)
+
\log p(A_1\mid S_1;\theta)
\nonumber\\
&\quad+
\sum_{t=2}^{T}\log p(S_t\mid S_{t-1};\theta)
+
\sum_{t=2}^{T}\log p(A_t\mid A_{t-1},S_t;\theta)
\nonumber\\
&\quad+
\sum_{t=1}^{T}\log p(O_t\mid S_t,A_t;\theta)
\label{eq:complete_data_loglik}
\end{align}

Substituting \eqref{eq:complete_data_loglik} into the auxiliary function
\eqref{eq:q_function_def} gives:
\begin{align}
Q(\theta,\theta^{\mathrm{old}})
&=
\mathbb{E}_{p(S_{1:T},A_{1:T}\mid O_{1:T},\theta^{\mathrm{old}})}
\left[
\log p(O_{1:T},S_{1:T},A_{1:T}\mid \theta)
\right]
\nonumber\\
&=
\mathbb{E}_{p(S_{1:T},A_{1:T}\mid O_{1:T},\theta^{\mathrm{old}})}
\Biggl[
\log p(S_1\mid \theta)
+
\log p(A_1\mid S_1;\theta)
\nonumber\\
&\qquad\qquad\qquad\qquad
+
\sum_{t=2}^{T}\log p(S_t\mid S_{t-1};\theta)
+
\sum_{t=2}^{T}\log p(A_t\mid A_{t-1},S_t;\theta)
\nonumber\\
&\qquad\qquad\qquad\qquad
+
\sum_{t=1}^{T}\log p(O_t\mid S_t,A_t;\theta)
\Biggr]
\label{eq:q_expanded_inside_expectation}
\end{align}

By linearity of expectation, this becomes
\begin{align}
Q(\theta,\theta^{\mathrm{old}})
&=
\mathbb{E}\!\left[\log p(S_1\mid \theta)\right]
+
\mathbb{E}\!\left[\log p(A_1\mid S_1;\theta)\right]
\nonumber\\
&\quad+
\sum_{t=2}^{T}
\mathbb{E}\!\left[\log p(S_t\mid S_{t-1};\theta)\right]
+
\sum_{t=2}^{T}
\mathbb{E}\!\left[\log p(A_t\mid A_{t-1},S_t;\theta)\right]
\nonumber\\
&\quad+
\sum_{t=1}^{T}
\mathbb{E}\!\left[\log p(O_t\mid S_t,A_t;\theta)\right],
\label{eq:q_linearity}
\end{align}
where all expectations are taken with respect to
\(p(S_{1:T},A_{1:T}\mid O_{1:T},\theta^{\mathrm{old}})\) \cite{Dempster1977,pml1Book}.

These expectations can be written in terms of the posterior marginals from the E-step. 
For the initial style term, define:
\begin{equation}
\gamma_1^{S}(s)
=
p(S_1=s\mid O_{1:T},\theta^{\mathrm{old}})
=
\sum_a \gamma_1(s,a)
\label{eq:gamma_style_init}
\end{equation}
Using \eqref{eq:gamma_joint_appendix}, \eqref{eq:xi_style_appendix}, and
\eqref{eq:xi_action_appendix}, the auxiliary function becomes:
\begin{align}
Q(\theta,\theta^{\mathrm{old}})
&=
\sum_s \gamma_1^{S}(s)\,\log p(S_1=s\mid \theta)
+
\sum_{s,a}\gamma_1(s,a)\,\log p(A_1=a\mid S_1=s;\theta)
\nonumber\\
&\quad+
\sum_{t=2}^{T}\sum_{s',s}
\xi_t^{S}(s',s)\,
\log p(S_t=s\mid S_{t-1}=s';\theta)
\nonumber\\
&\quad+
\sum_{t=2}^{T}\sum_{s,a',a}
\xi_t^{A}(s,a',a)\,
\log p(A_t=a\mid A_{t-1}=a',S_t=s;\theta)
\nonumber\\
&\quad+
\sum_{t=1}^{T}\sum_{s,a}
\gamma_t(s,a)\,
\log p(O_t\mid S_t=s,A_t=a;\theta)
\label{eq:q_fully_expanded}
\end{align}
This is the expected complete-data log-likelihood written explicitly in terms of the posterior quantities computed in the E-step \cite{Dempster1977,rabiner1989tutorial,murphy2002dynamic}.

For the emission-parameter derivations in this appendix, only the last term in
\eqref{eq:q_fully_expanded} is required. Therefore, the emission contribution is:
\begin{equation}
Q_{\mathrm{emit}}(\theta,\theta^{\mathrm{old}})
=
\sum_{t=1}^{T}\sum_{s,a}
\gamma_t(s,a)\,
\log p(O_t\mid S_t=s,A_t=a;\theta)
\label{eq:q_emit_only}
\end{equation}
The following sections specialize \eqref{eq:q_emit_only} to the hierarchical and \ac{PoE} emission formulations considered in this work.





\section{Hierarchical Emission Model}
\label{app:hierarchical_emission}

For the hierarchical emission formulation introduced in \autoref{sec:proposed_dbn_model}, 
the emission model is parameterized directly for each joint latent-state pair \((s,a)\). 
The corresponding M-step therefore follows the standard \ac{EM} pattern for latent-variable 
models, where each parameter block is updated using posterior weights \cite{Dempster1977,Bilmes1998}.

Using the decomposition $O_t=(x_t,b_t)$, the emission factorizes as:
\begin{equation}
p(O_t \mid s,a;\theta_{s,a})
=
p(x_t \mid s,a;\theta^{(x)}_{s,a})\,
p(b_t \mid s,a;\theta^{(b)}_{s,a}),
\end{equation}
with
\begin{equation}
\theta^{(x)}_{s,a} = (\mu_{s,a}, \Sigma_{s,a}),
\qquad
\theta^{(b)}_{s,a} = \{p_{s,a,j}\}_{j=1}^{B}
\end{equation}


For the continuous features, a Gaussian distribution with diagonal covariance is used:
\begin{equation}
p(x_t \mid s,a)
=
\mathcal{N}(x_t \mid \mu_{s,a},\Sigma_{s,a}),
\qquad
\Sigma_{s,a}
=
\mathrm{diag}\!\left(\sigma^2_{s,a,1},\dots,\sigma^2_{s,a,D_c}\right)
\end{equation}

For the binary features, a factorized Bernoulli distribution is used:
\begin{equation}
p(b_t \mid s,a)
=
\prod_{j=1}^{B}
(p_{s,a,j})^{b_t^{(j)}}(1-p_{s,a,j})^{1-b_t^{(j)}}
\end{equation}


Hence, for a fixed joint state pair \((s,a)\), the log-emission probability becomes:
\begin{align}
\log p(O_t \mid s,a;\theta_{s,a})
&=
\log p(x_t \mid s,a)
+
\log p(b_t \mid s,a)
\nonumber\\
&=
\log \mathcal{N}(x_t \mid \mu_{s,a},\Sigma_{s,a})
+
\sum_{j=1}^{B}
\left[
b_t^{(j)}\log p_{s,a,j}
+
(1-b_t^{(j)})\log(1-p_{s,a,j})
\right]
\label{eq:hier_log_emission}
\end{align}

Substituting \eqref{eq:hier_log_emission} into the emission contribution of the auxiliary function,
defined in \eqref{eq:q_emit_only}, gives:
\begin{align}
Q_{\mathrm{emit}}(\theta,\theta^{\mathrm{old}})
&=
\sum_{t=1}^{T}\sum_{s,a}
\gamma_t(s,a)\,
\log \mathcal{N}(x_t \mid \mu_{s,a},\Sigma_{s,a})
\nonumber\\
&\quad+
\sum_{t=1}^{T}\sum_{s,a}
\gamma_t(s,a)
\sum_{j=1}^{B}
\left[
b_t^{(j)}\log p_{s,a,j}
+
(1-b_t^{(j)})\log(1-p_{s,a,j})
\right]
\label{eq:hier_qemit_expanded}
\end{align}


For convenience, define the effective posterior count of the joint state pair \((s,a)\) as:
\begin{equation}
N_{s,a}
=
\sum_{t=1}^{T}\gamma_t(s,a)
\label{eq:hier_effective_count}
\end{equation}


\subsection{Parameter Updates}

The emission parameters are obtained by maximizing \eqref{eq:hier_qemit_expanded} with respect to 
the parameters of each joint state pair \((s,a)\). Because the parameters for one joint state do not 
appear in the emission terms of any other joint state, the optimization separates across \((s,a)\) 
and yields closed-form updates \cite{Dempster1977,Bilmes1998}.

\paragraph{Gaussian mean.}
For fixed covariance, the part of the objective that depends on \(\mu_{s,a}\) is:
\begin{equation}
-\frac{1}{2}
\sum_{t=1}^{T}
\gamma_t(s,a)
(x_t-\mu_{s,a})^\top
\Sigma_{s,a}^{-1}
(x_t-\mu_{s,a})
\end{equation}
Differentiating with respect to \(\mu_{s,a}\) and setting the derivative to zero yields:
\begin{equation}
\mu_{s,a}
=
\frac{\sum_{t=1}^{T}\gamma_t(s,a)\,x_t}
{\sum_{t=1}^{T}\gamma_t(s,a)}
=
\frac{1}{N_{s,a}}
\sum_{t=1}^{T}\gamma_t(s,a)\,x_t
\label{eq:hier_mean_update}
\end{equation}


\paragraph{Gaussian covariance.}
Assuming a diagonal covariance matrix, each variance parameter can be optimized independently. 
For dimension \(d\), the relevant term is
\begin{equation}
-\frac{1}{2}
\sum_{t=1}^{T}
\gamma_t(s,a)
\left[
\log \sigma^2_{s,a,d}
+
\frac{\bigl(x_t^{(d)}-\mu_{s,a}^{(d)}\bigr)^2}{\sigma^2_{s,a,d}}
\right].
\end{equation}
Differentiating with respect to \(\sigma^2_{s,a,d}\) and setting the result to zero gives:
\begin{equation}
\sigma^2_{s,a,d}
=
\frac{1}{N_{s,a}}
\sum_{t=1}^{T}
\gamma_t(s,a)\,
\bigl(x_t^{(d)}-\mu_{s,a}^{(d)}\bigr)^2
\label{eq:hier_var_update}
\end{equation}
Hence,
\begin{equation}
\Sigma_{s,a}
=
\mathrm{diag}\!\left(\sigma^2_{s,a,1},\dots,\sigma^2_{s,a,D_c}\right).
\end{equation}

\paragraph{Bernoulli parameters.}
For each binary feature \(j\), the Bernoulli part of the objective is
\begin{equation}
\sum_{t=1}^{T}
\gamma_t(s,a)
\left[
b_t^{(j)}\log p_{s,a,j}
+
(1-b_t^{(j)})\log(1-p_{s,a,j})
\right].
\end{equation}
Differentiating with respect to \(p_{s,a,j}\) and setting this derivative to zero gives:
\begin{equation}
p_{s,a,j}
=
\frac{\sum_{t=1}^{T}\gamma_t(s,a)\,b_t^{(j)}}
{\sum_{t=1}^{T}\gamma_t(s,a)}
=
\frac{1}{N_{s,a}}
\sum_{t=1}^{T}\gamma_t(s,a)\,b_t^{(j)}
\label{eq:hier_bernoulli_update}
\end{equation}

Thus, the hierarchical emission model admits closed-form M-step updates, 
where each parameter is estimated as a posterior-weighted maximum-likelihood quantity for the corresponding joint latent-state pair \cite{Dempster1977,Bilmes1998}.




\section{Product-of-Experts Emission Model}
\label{app:poe_emission}

For the \ac{PoE} emission formulation introduced in \autoref{sec:proposed_dbn_model}, 
the emission distribution for the joint latent-state pair \((s,a)\) is constructed by 
multiplying a style-dependent expert and an action-dependent expert, followed by normalization:
\begin{equation}
p(O_t \mid S_t=s,A_t=a)
=
\frac{p_s(O_t \mid s)\,p_a(O_t \mid a)}{Z_{s,a}}
\label{eq:poe_emission_def}
\end{equation}
This is the standard product-of-experts construction, in which several simpler expert models 
are combined multiplicatively and then renormalized to obtain a valid probability distribution 
\cite{Hinton2002}.

Here, \(Z_{s,a}\) is the normalization constant that ensures that
\eqref{eq:poe_emission_def} defines a valid probability distribution. Using the
continuous-binary decomposition of the observation, it can be written more explicitly as:
\begin{equation}
Z_{s,a}
=
\sum_{b \in \{0,1\}^{B}}
\int
p_s(x,b \mid s)\,p_a(x,b \mid a)\,dx
\end{equation}

Using the decomposition \(O_t=(x_t,b_t)\), both experts factorize into continuous and binary parts:
\begin{align}
p_s(O_t \mid s)
&=
p_s(x_t \mid s)\,p_s(b_t \mid s),
\\
p_a(O_t \mid a)
&=
p_a(x_t \mid a)\,p_a(b_t \mid a)
\end{align}

The corresponding parameter sets are written as:
\begin{equation}
\theta_s=(\mu_s,\Sigma_s,\{p_j^s\}_{j=1}^{B}),
\qquad
\theta_a=(\mu_a,\Sigma_a,\{p_j^a\}_{j=1}^{B})
\end{equation}

For the continuous features, Gaussian experts with diagonal covariance are used:
\begin{align}
p_s(x_t \mid s)
&=
\mathcal{N}(x_t \mid \mu_s,\Sigma_s),
\qquad
\Sigma_s=\mathrm{diag}\!\left(\sigma^2_{s,1},\dots,\sigma^2_{s,D_c}\right),
\\
p_a(x_t \mid a)
&=
\mathcal{N}(x_t \mid \mu_a,\Sigma_a),
\qquad
\Sigma_a=\mathrm{diag}\!\left(\sigma^2_{a,1},\dots,\sigma^2_{a,D_c}\right)
\end{align}

For the binary features, factorized Bernoulli experts are used:
\begin{align}
p_s(b_t \mid s)
&=
\prod_{j=1}^{B}
(p_j^s)^{b_t^{(j)}}(1-p_j^s)^{1-b_t^{(j)}},
\\
p_a(b_t \mid a)
&=
\prod_{j=1}^{B}
(p_j^a)^{b_t^{(j)}}(1-p_j^a)^{1-b_t^{(j)}}
\end{align}

Because the experts factorize into continuous and binary parts, the normalization constant also factorizes:
\begin{align}
Z_{s,a}
&=
\sum_{b \in \{0,1\}^{B}}
\int
p_s(x,b \mid s)\,p_a(x,b \mid a)\,dx
\nonumber\\
&=
\left(
\int p_s(x \mid s)\,p_a(x \mid a)\,dx
\right)
\left(
\sum_{b \in \{0,1\}^{B}}
p_s(b \mid s)\,p_a(b \mid a)
\right)
\nonumber\\
&=
Z_{s,a}^{(x)}\,Z_{s,a}^{(b)}
\end{align}

The continuous normalization term is therefore:
\begin{equation}
Z_{s,a}^{(x)}
=
\int
\mathcal{N}(x \mid \mu_s,\Sigma_s)\,
\mathcal{N}(x \mid \mu_a,\Sigma_a)\,dx
\end{equation}
Using the standard Gaussian product identity, this becomes
\begin{equation}
Z_{s,a}^{(x)}
=
\mathcal{N}(\mu_s \mid \mu_a,\Sigma_s+\Sigma_a)
\label{eq:poe_zx_closed}
\end{equation}
Hence,
\begin{align}
\log Z_{s,a}^{(x)}
&=
-\frac{D_c}{2}\log(2\pi)
-\frac{1}{2}\log |\Sigma_s+\Sigma_a|
\nonumber\\
&\quad
-\frac{1}{2}
(\mu_s-\mu_a)^\top
(\Sigma_s+\Sigma_a)^{-1}
(\mu_s-\mu_a)
\end{align}

For the binary part,
\begin{align}
Z_{s,a}^{(b)}
&=
\sum_{b \in \{0,1\}^{B}}
\prod_{j=1}^{B}
(p_j^s)^{b^{(j)}}(1-p_j^s)^{1-b^{(j)}}
(p_j^a)^{b^{(j)}}(1-p_j^a)^{1-b^{(j)}}
\end{align}
Since the Bernoulli experts factorize across dimensions, the sum also factorizes:
\begin{equation}
Z_{s,a}^{(b)}
=
\prod_{j=1}^{B}
\left[
p_j^s p_j^a + (1-p_j^s)(1-p_j^a)
\right]
\label{eq:poe_zb_closed}
\end{equation}
Taking the logarithm yields
\begin{equation}
\log Z_{s,a}^{(b)}
=
\sum_{j=1}^{B}
\log
\left[
p_j^s p_j^a + (1-p_j^s)(1-p_j^a)
\right]
\end{equation}

Combining the two experts and the normalization term, the log-emission probability becomes:
\begin{align}
\log p(O_t \mid s,a)
&=
\log p_s(x_t \mid s)
+
\log p_a(x_t \mid a)
+
\log p_s(b_t \mid s)
+
\log p_a(b_t \mid a)
\nonumber\\
&\quad
-
\log Z_{s,a}^{(x)}
-
\log Z_{s,a}^{(b)}
\label{eq:poe_log_emission}
\end{align}


Substituting \eqref{eq:poe_log_emission} into the emission contribution of the auxiliary function,
defined in \eqref{eq:q_emit_only}, gives:
\begin{align}
Q_{\mathrm{emit}}(\theta,\theta^{\mathrm{old}})
&=
\sum_{t=1}^{T}\sum_{s,a}
\gamma_t(s,a)\,
\log \mathcal{N}(x_t \mid \mu_s,\Sigma_s)
\nonumber\\
&\quad+
\sum_{t=1}^{T}\sum_{s,a}
\gamma_t(s,a)\,
\log \mathcal{N}(x_t \mid \mu_a,\Sigma_a)
\nonumber\\
&\quad+
\sum_{t=1}^{T}\sum_{s,a}
\gamma_t(s,a)
\sum_{j=1}^{B}
\left[
b_t^{(j)}\log p_j^s + (1-b_t^{(j)})\log(1-p_j^s)
\right]
\nonumber\\
&\quad+
\sum_{t=1}^{T}\sum_{s,a}
\gamma_t(s,a)
\sum_{j=1}^{B}
\left[
b_t^{(j)}\log p_j^a + (1-b_t^{(j)})\log(1-p_j^a)
\right]
\nonumber\\
&\quad-
\sum_{t=1}^{T}\sum_{s,a}
\gamma_t(s,a)\,\log Z_{s,a}^{(x)}
-
\sum_{t=1}^{T}\sum_{s,a}
\gamma_t(s,a)\,\log Z_{s,a}^{(b)}
\label{eq:poe_qemit_expanded}
\end{align}

\subsection{Implication for the M-Step}

In contrast to the hierarchical emission model, the \ac{PoE} objective does not separate into 
independent parameter blocks for each joint state pair. The reason is that the normalization terms 
\(\log Z_{s,a}^{(x)}\) and \(\log Z_{s,a}^{(b)}\) depend on the same expert parameters that also 
appear in the data-fit terms. As a result, the style and action experts are coupled through the 
partition function \cite{Hinton2002}.

More specifically, any update of a style-expert parameter affects not only the style likelihood term, 
but also all normalization terms involving that style state across the action experts. The same holds 
symmetrically for the action-expert parameters. Therefore, the emission M-step no longer reduces to a 
set of independent closed-form weighted maximum-likelihood updates.

At the theoretical level, this means that the \ac{PoE} emission parameters are optimized by increasing 
$
Q_{\mathrm{emit}}(\theta,\theta^{\mathrm{old}})
$ 
with respect to the expert parameters using a generalized \ac{EM} step, rather than by solving 
closed-form maximization equations for each parameter separately \cite{Dempster1977,Wu1983}.

\subsection{Gradient Form}

Let \(\vartheta_s\) denote a parameter of the style expert and \(\vartheta_a\) a parameter of the 
action expert. Differentiating \eqref{eq:poe_qemit_expanded} gives the generic gradient structure:
\begin{equation}
\frac{\partial Q_{\mathrm{emit}}}{\partial \vartheta_s}
=
\sum_{t=1}^{T}\sum_{a}
\gamma_t(s,a)
\left[
\frac{\partial}{\partial \vartheta_s}\log p_s(O_t \mid s)
-
\frac{\partial}{\partial \vartheta_s}\log Z_{s,a}
\right],
\label{eq:poe_style_grad}
\end{equation}
and similarly
\begin{equation}
\frac{\partial Q_{\mathrm{emit}}}{\partial \vartheta_a}
=
\sum_{t=1}^{T}\sum_{s}
\gamma_t(s,a)
\left[
\frac{\partial}{\partial \vartheta_a}\log p_a(O_t \mid a)
-
\frac{\partial}{\partial \vartheta_a}\log Z_{s,a}
\right]
\label{eq:poe_action_grad}
\end{equation}

These expressions show that each gradient contains two competing contributions:
a data-fit term from the expert likelihood and a normalization term arising from the partition 
function. The second term is specific to the \ac{PoE} formulation and is the reason why the 
emission M-step does not admit the same closed-form updates as in the hierarchical case 
\cite{Hinton2002}.